\let\negmedspace\undefined
\let\negthickspace\undefined
\documentclass[journal]{IEEEtran}
\usepackage[a5paper, margin=10mm, onecolumn]{geometry}
%\usepackage{lmodern} % Ensure lmodern is loaded for pdflatex
\usepackage{tfrupee} % Include tfrupee package

\setlength{\headheight}{1cm} % Set the height of the header box
\setlength{\headsep}{0mm}     % Set the distance between the header box and the top of the text

\usepackage{gvv-book}
\usepackage{gvv}
\usepackage{cite}
\usepackage{amsmath,amssymb,amsfonts,amsthm}
\usepackage{algorithmic}
\usepackage{graphicx}
\usepackage{textcomp}
\usepackage{xcolor}
\usepackage{txfonts}
\usepackage{listings}
\usepackage{enumitem}
\usepackage{mathtools}
\usepackage{gensymb}
\usepackage[breaklinks=true]{hyperref}
\usepackage{tkz-euclide} 
\usepackage{listings}
% \usepackage{gvv}                                        
\def\inputGnumericTable{}                   
\usepackage[latin1]{inputenc}                                
\usepackage{color}                                            
\usepackage{array}                                            
\usepackage{longtable}                                       
\usepackage{calc}                                             
\usepackage{multirow}                                         
\usepackage{hhline}                                           
\usepackage{ifthen}                                           
\usepackage{lscape}
\usepackage{circuitikz}
\usepackage{comment}
\tikzstyle{block} = [rectangle, draw, fill=blue!20, 
text width=4em, text centered, rounded corners, minimum height=3em]
\tikzstyle{sum} = [draw, fill=blue!10, circle, minimum size=1cm, node distance=1.5cm]
\tikzstyle{input} = [coordinate]
\tikzstyle{output} = [coordinate]
%--------------------------

\begin{document}
	
	\bibliographystyle{IEEEtran}
	\vspace{3cm}
	
	\title{4.13.46}
	\author{EE25BTECH11042 - Nipun Dasari}
	\maketitle
	
	\renewcommand{\thefigure}{\theenumi}
	\renewcommand{\thetable}{\theenumi}
	\setlength{\intextsep}{10pt}
	
	\numberwithin{equation}{enumi}
	\numberwithin{figure}{enumi}
	\renewcommand{\thetable}{\theenumi}
	
	\textbf{Question}:\\
	A straight line L is perpendicular to the line $5x - y + 1 = 0$. The area of the triangle formed by L and the coordinate axes is 5. Find the equation of the line L.
	
	\solution
	
	Let the given line be $L_1$ and the required perpendicular line be $L$.
	The equation of the given line $L_1$ is $5x - y = -1$. Using the matrix form $\vec{n}^\top\vec{x} = c$, its parameters are:
	\begin{align}
		\vec{n}_1 = \begin{myvec}{5 \\ -1}\end{myvec}, \quad \vec{x} = \begin{myvec}{x \\ y}\end{myvec}, \quad c_1 = -1
	\end{align}
	Let the required line $L$ have the equation $\vec{n}_2^\top\vec{x} = c_2$. Since the lines are perpendicular, the dot product of their normal vectors must be zero, so $\vec{n}_1^\top\vec{n}_2 = 0$.
	
	Let's assume the normal vector for line L is $\vec{n}_2 = \begin{myvec}{1 \\ m}\end{myvec}$. We can solve for $m$:
	\begin{align}
		\vec{n}_1^\top\vec{n}_2 = \begin{myvec}{5 & -1}\end{myvec}\begin{myvec}{1 \\ m}\end{myvec} = 5 - m &= 0 \nonumber \\
		\implies m &= 5
	\end{align}
	Therefore, the normal vector for line L is $\vec{n}_2 = \begin{myvec}{1 \\ 5}\end{myvec}$, and its equation is $x + 5y = c_2$.
	
	The line L forms a triangle with the coordinate axes. We find the intercepts. For the x-intercept, set $y=0$, which gives $x = c_2$. The point is $(c_2, 0)$. For the y-intercept, set $x=0$, which gives $5y = c_2$, so $y = c_2/5$. The point is $(0, c_2/5)$.
	
	The area of this triangle is given as 5. Using the area formula:
	\begin{align}
		\text{Area} = \frac{1}{2} |(\text{x-intercept}) \cdot (\text{y-intercept})| &= 5 \nonumber \\
		\frac{1}{2} |(c_2)(c_2/5)| &= 5 \nonumber \\
		|\frac{c_2^2}{10}| &= 5 \nonumber \\
		c_2^2 &= 50 \implies c_2 = \pm\sqrt{50} = \pm 5\sqrt{2}
	\end{align}
	
	There are two possible lines that satisfy the given conditions. Substituting the values of $c_2$ into the matrix equation $\vec{n}_2^\top\vec{x} = c_2$, the possible equations for line L are:
	\begin{align}
		\begin{myvec}{1 & 5}\end{myvec}\begin{myvec}{x \\ y}\end{myvec} &= 5\sqrt{2} \\
		\begin{myvec}{1 & 5}\end{myvec}\begin{myvec}{x \\ y}\end{myvec} &= -5\sqrt{2}
	\end{align}
	
\end{document}